\documentclass[10pt,twocolumn]{article}
\usepackage[margin=0.75in,top=0.90in,bottom=0.90in]{geometry}

\makeatletter
\def\ps@acmprobe{%
  \def\@oddhead{\small SurvEval-Probe: Two-Column Ranking Evaluation\hfill Conference '26, November 2026}%
  \def\@evenhead{\small Conference '26, November 2026\hfill Rivera and Vance}%
  \def\@oddfoot{\hfill\small\thepage\hfill}%
  \def\@evenfoot{\hfill\small\thepage\hfill}%
}
\makeatother
\pagestyle{acmprobe}

\title{\textbf{Survival-Weighted Listwise Ranking Evaluation}}
\author{Alice Rivera \and Bernard Vance}
\date{November 2026}

\begin{document}
\maketitle
\thispagestyle{acmprobe}

\section*{ABSTRACT}
We evaluate geometric passage segmentation on two-column conference articles featuring all-caps section headings, spanning benchmark tables, and hanging-indented numbered reference lists. Across 420 evaluation queries, our method preserves bibliography items as body passages.

\section*{1 INTRODUCTION}
Modern conference templates such as ACM sigconf typeset top-level headings in uppercase and format bibliographies as two-column hanging-indented lists with bracketed numeric labels. Naive table detectors merge left and right bibliography columns across the central gutter.

\section*{2 METHODOLOGY}
Each page is partitioned into margin furniture, genuine cross-column floats, and two-column prose streams. Only multi-column grids or captioned floats that bridge across the centerline are isolated as spanning figures or tables.

\begin{table*}[t]
\centering
\caption{Cross-column evaluation of listwise ranking metrics across two benchmark datasets.}
\vspace{3pt}
\begin{tabular}{lcccc}
\hline
Model Variant & Dataset A NDCG & Dataset A uAUC & Dataset B NDCG & Latency \\
\hline
Pointwise Baseline & 0.7965 & 0.6781 & 0.7022 & 6.4 ms \\
Permutation Ranker & 0.8312 & 0.7091 & 0.7281 & 11.2 ms \\
Contextual Evaluator & 0.8450 & 0.7196 & 0.7371 & 12.9 ms \\
Proposed Listwise Model & 0.8523 & 0.7236 & 0.7450 & 14.1 ms \\
\hline
\end{tabular}
\end{table*}

\section*{3 EXPERIMENTS}
Table 1 reports ranking performance across both benchmark datasets. Separating true cross-column floats from single-column prose ensures that section boundaries and citation entries remain searchable.

\section*{4 CONCLUSION}
Distinguishing hanging-indented reference lists from tabular grids preserves both the uppercase REFERENCES heading and individual bibliographic entries in their proper reading columns.

\clearpage
\section*{REFERENCES}
{\small
\begin{list}{[\arabic{enumi}]}{%
  \usecounter{enumi}%
  \setlength{\leftmargin}{2.2em}%
  \setlength{\labelwidth}{1.3em}%
  \setlength{\labelsep}{0.9em}%
  \setlength{\itemsep}{2pt}%
  \setlength{\parsep}{0pt}%
}
\item Qingyao Ai, Keping Bi, Jiafeng Guo, and W. Bruce Croft. 2018. Learning a deep listwise context model for ranking refinement. In Proceedings of the 41st International ACM SIGIR Conference. 135--144.
\item Irwan Bello, Sayali Kulkarni, Sagar Jain, Craig Boutilier, and Ofer Meshi. 2018. Seq2Slate: Re-ranking and slate optimization with pointer networks. arXiv preprint arXiv:1810.02019 (2018).
\item Yukuo Cen, Jianwei Zhang, Xu Zou, Chang Zhou, and Jie Tang. 2020. Controllable multi-interest framework for recommendation. In Proceedings of the 26th ACM SIGKDD Conference. 2942--2951.
\item Paul Covington, Jay Adams, and Emre Sargin. 2016. Deep neural networks for video recommendations. In Proceedings of the 10th ACM Conference on Recommender Systems. 191--198.
\item Alex Kendall, Yarin Gal, and Roberto Cipolla. 2018. Multi-task learning using uncertainty to weigh losses for scene geometry. In Proceedings of the IEEE CVPR Conference. 7482--7491.
\newpage
\item Changhua Pei, Yi Zhang, Yongfeng Zhang, Fei Sun, Xiao Lin, and Wenwu Ou. 2019. Personalized re-ranking for recommendation. In Proceedings of the 13th ACM Conference on Recommender Systems. 3--11.
\item Yuxin Ren, Qiya Yang, Yichun Wu, Wei Xu, and Zhiqiang Zhang. 2024. Non-autoregressive generative models for reranking recommendation. In Proceedings of the 30th ACM SIGKDD Conference. 5625--5634.
\item Xiaowen Shi, Fan Yang, Ze Wang, Xiaoxu Wu, and Dong Wang. 2023. Permutation-level interest-based end-to-end re-ranking framework. In Proceedings of the 29th ACM SIGKDD Conference. 4823--4831.
\item Liang Pang, Jun Xu, Qingyao Ai, Yanyan Lan, and Jirong Wen. 2020. SetRank: Learning a permutation-invariant ranking model for information retrieval. In Proceedings of the 43rd ACM SIGIR Conference. 499--508.
\item Tao Zhuang, Wenwu Ou, and Zhirong Wang. 2018. Globally optimized mutual influence aware ranking in search. arXiv preprint arXiv:1805.08524 (2018).
\end{list}
}
\end{document}
